% ===========================================================================
%  Nonlinear Algebra -- fill-in notes on Chapters 1-2
%  (Michalek & Sturmfels, "Invitation to Nonlinear Algebra", AMS GSM 211)
%  woven on LOOM.   compile:  latexmk -xelatex main   (or xelatex twice)
%
%  HOW THESE NOTES WORK
%   - Read the strands and the stated theorems (passive intake).
%   - Fill the blanks \fillin, the \TODO proof gaps, and the "your turn" boxes
%     (active recall). The book is your answer key.
% ===========================================================================
\documentclass[cjk]{loom}   % [cjk] only for the pen name; all prose is English

% --- pedagogy tools (\block \trigger \TODO \fillin \recall yourturn \workspace)
%     and booktabs/tabularx/L now come built into loom.cls v0.2 ---------------

% --- math shorthands --------------------------------------------------------
\newcommand{\R}{\mathbb{R}}
\newcommand{\C}{\mathbb{C}}
\newcommand{\Q}{\mathbb{Q}}
\newcommand{\Z}{\mathbb{Z}}
\newcommand{\N}{\mathbb{N}}
\newcommand{\PP}{\mathbb{P}}
\newcommand{\V}{\mathcal{V}}        % variety  V(I)
\newcommand{\Ideal}{\mathcal{I}}    % vanishing ideal  I(W)
\newcommand{\gen}[1]{\langle #1\rangle}
\newcommand{\xb}{\mathbf{x}}
\newcommand{\inn}{\operatorname{in}}
\DeclareMathOperator{\Spec}{Spec}
\DeclareMathOperator{\rad}{rad}
\DeclareMathOperator{\HS}{HS}
\DeclareMathOperator{\HP}{HP}
\DeclareMathOperator{\rank}{rank}

\runningthread{nonlinear algebra · ch 1--2}

\begin{document}

\loomcover
  {Nonlinear Algebra}
  {Fill-in notes · Chapters 1--2 · Michałek \& Sturmfels}
  {北极甜虾 (剑心犹在！)}
  {\today}

% ===========================================================================
\section*{The dictionary, and how to read these notes}
\addcontentsline{toc}{section}{The dictionary · how to read}
\warmth{0}

\begin{strand}
The whole of Chapters 1--2 is \emph{one} idea, stated twice by Sophie Germain:
\keyword{``Algebra is but written geometry''} (Ch.\,1) and
\keyword{``Geometry is but drawn algebra''} (Ch.\,2). There is a dictionary
between the \emph{algebra} side (ideals in the polynomial ring $K[\xb]$) and the
\emph{geometry} side (varieties in $K^n$ or $\PP^n$). Learn the dictionary and
the two chapters become one.
\end{strand}

\block{How to read these notes (passive \emph{and} active)}
\begin{itemize}[leftmargin=1.3em,itemsep=1pt]
  \item \keyword{Read} the strands and the stated theorems. That is the passive intake.
  \item \keyword{Fill} the blanks \fillin[1.2cm], the \TODO{the missing step} gaps inside proofs,
        and the \emph{Your turn} boxes. That is the active recall. The book is your answer key.
  \item Bump each section's opening \texttt{\textbackslash warmth\{0\}} to $0$--$5$ once you can rebuild it cold.
\end{itemize}

\block{The master dictionary (fill it in as you go)}
\noindent\begin{tabularx}{\linewidth}{@{}L L@{}}
\toprule
{\color{indigo}Algebra: ideals in $K[\xb]$} & {\color{madder}Geometry: varieties in $K^n,\PP^n$}\\
\midrule
ideal $I$                       & variety $\V(I)$ \;(§\ref{sec:affine})\\
vanishing ideal $\Ideal(W)$ (radical) & variety $W=\V(\Ideal(W))$\\
$I\subseteq J$                  & $\V(I)\;\fillin[1cm]\;\V(J)$ \;(which way?)\\
sum $I+J$                       & $\fillin[2cm]$ of $\V(I),\V(J)$\\
intersection $I\cap J$          & $\fillin[2cm]$ of $\V(I),\V(J)$\\
prime ideal                     & \fillin[3cm] variety \;(§\ref{sec:affine})\\
maximal ideal $\gen{x_1-p_1,\dots,x_n-p_n}$ & a single point $\fillin[1.5cm]$\\
radical ideal                   & ($\Ideal\!=\!\V$ inverse, up to radical: Nullstellensatz, Ch.\,6)\\
quotient $K[\xb]/\Ideal(W)$     & coordinate ring $K[W]$ = functions on $W$\\
dimension of $I$ (§\ref{sec:dimdeg}) & dimension of $\V(I)$\\
degree of $I$                   & \# points of a general linear slice\\
$\gen{f,\partial f/\partial x_1,\dots}$ has no zeros & $\V(f)$ is \fillin[2cm]\\
homogeneous ideal               & projective variety \;(§\ref{sec:proj})\\
primary decomposition (Ch.\,3)  & decomposition into irreducibles\\
\bottomrule
\end{tabularx}

\begin{strand}
The map $\V$ sends algebra to geometry; the map $\Ideal$ sends geometry back to
algebra. They are \emph{inclusion-reversing}, and they are almost inverse to each
other. ``Almost'' is the content of the Nullstellensatz (Chapter 6); for now,
$\Ideal(\V(I))=\rad(I)$ is the slogan to carry.
\end{strand}

\newpage
{\small\tableofcontents}

\clearpage
% ===========================================================================
% ===========================================================================
\section{Ideals}\label{sec:ideals}
\warmth{0}\quad\whisper{The algebra side. An ideal is the algebraic shadow of a geometric shape.}

\begin{strand}
The objects of nonlinear algebra live in the \keyword{polynomial ring}
$K[\xb]=K[x_1,\dots,x_n]$ over a field $K$ (favorite: $K=\Q$; also $\R,\C,\overline K,
\mathbb F_q$). A polynomial is a finite $K$-combination of monomials $\xb^{\mathbf a}
=x_1^{a_1}\cdots x_n^{a_n}$, and its \keyword{degree} is $\max\{|\mathbf a|:c_{\mathbf a}\ne0\}$.
The hero of the book is one cubic surface; meet it now, it returns everywhere.
\end{strand}

\begin{example}[the running cubic, $n=3$]\label{ex:cayley}
\[
  f \;=\; \det\!\begin{pmatrix}1&x&y\\ x&1&z\\ y&z&1\end{pmatrix}
    \;=\; 2xyz-x^2-y^2-z^2+1 .
\]
$\V(f)\subset\R^3$ is a cubic surface (Cayley's nodal cubic). Its four \keyword{singular points}
are the common zeros of $f$ and its three partials $\partial f/\partial x=2yz-2x$, etc.
\end{example}
\begin{yourturn}
Solve $2yz-2x=2xz-2y=2xy-2z=0$ together with $f=0$. Check the four singular points are
$(1,1,1),(1,-1,-1),(-1,1,-1),(-1,-1,1)$. \recall{Why must a singular point also satisfy $f=0$ here? (Euler's identity for homogeneous-ish $f$ -- or just check.)}
\workspace[3]
\end{yourturn}

\block{What an ideal is}

\begin{definition}[ideal, Def.\ 1.1]\label{def:ideal}
All rings here are commutative with $1$. A nonempty subset $I\subseteq R$ is an \keyword{ideal} if
\begin{enumerate}[label=(\alph*),leftmargin=2em,itemsep=0pt]
  \item $f\in R$ and $g\in I$ \;$\Rightarrow$\; $\fillin[2.5cm]\in I$ \quad(absorbs multiplication);
  \item $f,g\in I$ \;$\Rightarrow$\; $\fillin[2cm]\in I$ \quad(closed under addition).
\end{enumerate}
Equivalently, $I=\ker\phi$ for some ring homomorphism $\phi:R\to S$.
\end{definition}

\recall{Why is $\{$even integers$\}$ an ideal of $\Z$? Which homomorphism has it as kernel?}%
Given $\mathcal F\subseteq R$, the smallest ideal containing it is $\gen{\mathcal F}$, the ideal
\keyword{generated} by $\mathcal F$ (all polynomial-coefficient combinations). The quotient $R/I$
is a ring (Prop.\ 1.2), built from congruence classes $f+I$.

\begin{proposition}[four operations, Prop.\ 1.3]\label{prop:ops}
If $I,J\subseteq R$ are ideals, so are the sum $I+J$, the intersection $I\cap J$, the product
$IJ=\gen{fg:f\in I,g\in J}$, and the colon (quotient) $(I:J)=\{f\in R:fJ\subseteq I\}$.
\end{proposition}
\begin{proof}
We check $(I:J)$ is an ideal. For (a): $f\in R,\ g\in(I:J)\Rightarrow (fg)J=f(gJ)\subseteq fI\subseteq I$.
For (b): $f,g\in(I:J)\Rightarrow (f+g)J\subseteq fJ+gJ\subseteq I+I=I$. \TODO{do the same two checks for $I\cap J$}
\end{proof}

\begin{yourturn}
($n=1$, Example 1.4.) In $\Q[x]$ let $I=\gen{(x+2)^3}$ and $J=\gen{(x-1)(x+2)}$. Compute, as
principal ideals (factor first!):
\[
  \begin{aligned}
    I\cap J&=\fillin[2.6cm], & IJ&=\fillin[2.6cm],\\[2pt]
    I+J&=\fillin[2.6cm],     & (I:J)&=\fillin[2.6cm].
  \end{aligned}
\]
Moral: \emph{arithmetic in $\Q[x]$ is just like arithmetic in $\Z$} ($\cap=\mathrm{lcm}$, $+=\gcd$).
\end{yourturn}

\block{Four kinds of ideal, read off the quotient}

\begin{definition}[the property dictionary]\label{def:primes}
For an ideal $I\subseteq R$ and its quotient $R/I$:
\end{definition}
\noindent\begin{tabularx}{\linewidth}{@{}l l L@{}}
\toprule
{\color{madder}$I$ is\dots} & {\color{indigo}definition} & \multicolumn{1}{l}{\color{indigo}\dots iff $R/I$ is}\\
\midrule
maximal  & no proper ideal contains $I$           & a \fillin[2cm]\\
prime    & $fg\in I\Rightarrow f\in I$ or $g\in I$ & an \fillin[3cm]\\
radical  & $f^s\in I\Rightarrow f\in I$            & free of \fillin[2.5cm] elements\\
primary  & $fg\in I,\ g\notin I\Rightarrow f^s\in I$ & one where zero divisors are nilpotent\\
\bottomrule
\end{tabularx}

\begin{proposition}[implications, Prop.\ 1.6]\label{prop:impl}
$I$ maximal $\Rightarrow$ prime $\Rightarrow$ radical, and prime $\Rightarrow$ primary.
None reverses; radical $+$ primary $=$ prime; any intersection of primes is radical.
\end{proposition}
\begin{proof}
Maximal $\Rightarrow$ prime since a field has no zero divisors. Prime $\Rightarrow$ radical:
take $g=f^{s-1}$ and induct on $s$. \TODO{prove radical $+$ primary $\Rightarrow$ prime (assume $fg\in I$, $f\notin I$; use primary then radical)}
\end{proof}

\recall{Match the non-reversing counterexamples in $\R[x,y]$: $\gen{x^2}$, $\gen{x(x-1)}$, $\gen{x}$ to ``primary not radical'', ``radical not primary'', ``prime not maximal''.}

\begin{strand}
Back to the cubic (Example~\ref{ex:cayley}). Its singular locus is $\Ideal=\gen{yz-x,\,xz-y,\,xy-z}$,
a \emph{radical} ideal that factors as an intersection of \keyword{five} maximal ideals; by the
Chinese Remainder Theorem $\R[x,y,z]/\Ideal\cong\R^5$. The number five is the five complex points
of $\V(\Ideal)$, four of them the real nodes. Algebra (an intersection of maximal ideals) has just
\emph{counted geometry} (a finite set of points). That is the dictionary firing for the first time.
\end{strand}

\loose{Exercise 18 (book): identify the maximal, prime, radical, primary ideals of $R=\Z$. Same four words, older ring.}

% ===========================================================================
\section{Gröbner Bases}\label{sec:groebner}
\warmth{0}\quad\whisper{The engine. Gröbner is to ideals what Gaussian elimination is to matrices.}

\begin{strand}
An ideal has many generating sets and no canonical basis. \keyword{Gröbner bases} fix this:
they are the right generating set once you pick a \emph{monomial order}. Think of them as the
Euclidean algorithm for $n\ge2$ variables, or Gaussian elimination for degree $\ge2$. The book's
creed: \keyword{nonlinear algebra is the next step after linear algebra}.
\end{strand}

\block{Monomial orders and initial terms}

\begin{definition}[monomial order, Def.\ 1.10]\label{def:order}
A total order $\prec$ on $\N^n$ (monomials) is a \keyword{monomial order} if
$(0,\dots,0)\preceq\mathbf a$ for all $\mathbf a$, and $\mathbf a\preceq\mathbf b\Rightarrow
\mathbf a+\mathbf c\preceq\mathbf b+\mathbf c$. Three standards (all with $x_1\succ\cdots\succ x_n$):
\emph{lex}, \emph{deglex} (degree first, then lex), \emph{degrevlex} (degree first, then
\emph{smallest} last variable wins). Any monomial order refines divisibility:
$\xb^{\mathbf a}\mid\xb^{\mathbf b}\Rightarrow\mathbf a\preceq\mathbf b$.
\end{definition}

The \keyword{initial monomial} $\inn_\prec(f)$ is the $\prec$-largest monomial of $f$.

\begin{yourturn}
Let $f=x^2+xz^2+y^3$ with $x\succ y\succ z$. Fill the three initial monomials:
\[
  \inn_{\mathrm{lex}}(f)=\fillin[1.6cm],\qquad
  \inn_{\mathrm{deglex}}(f)=\fillin[1.6cm],\qquad
  \inn_{\mathrm{degrevlex}}(f)=\fillin[1.6cm].
\]
\recall{For degrevlex, why does $y^3$ beat $xz^2$? Both have degree 3 -- compare the \emph{last} variable.}
\end{yourturn}

\block{Finiteness: Dickson and the Basis Theorem}

\begin{theorem}[Dickson's Lemma, Thm.\ 1.8]\label{thm:dickson}
Any infinite subset of $\N^n$ contains a pair $\{\mathbf a,\mathbf b\}$ with $\mathbf a\le\mathbf b$
(coordinatewise). Equivalently: any nonempty $\mathcal M\subseteq\N^n$ has a \emph{finite} set of
coordinatewise-minimal elements.
\end{theorem}
\begin{proof}
Induction on $n$. For $n=1$, any two naturals are comparable. \TODO{induction step: slice $\mathcal M$ by the last coordinate $i$, set $\mathcal M_i\subseteq\N^{n-1}$; if some $\mathcal M_i$ is infinite use IH, else \dots}
\end{proof}

For an ideal $I$, the \keyword{initial ideal} is $\inn_\prec(I)=\gen{\inn_\prec(f):f\in I\setminus\{0\}}$,
a monomial ideal. A finite $\mathcal G\subseteq I$ with $\inn_\prec(I)=\gen{\inn_\prec(g):g\in\mathcal G}$
is a \keyword{Gröbner basis} (it exists, by Dickson).

\begin{theorem}[a Gröbner basis generates, Thm.\ 1.13]\label{thm:gb-gen}
If $\mathcal G$ is a Gröbner basis of $I$, then $I=\gen{\mathcal G}$.
\end{theorem}
\begin{proof}
Suppose not. Pick $f\in I\setminus\gen{\mathcal G}$ with $\inn_\prec(f)=\xb^{\mathbf b}$ minimal.
Some $g\in\mathcal G$ has $\inn_\prec(g)\mid\xb^{\mathbf b}$, say $\xb^{\mathbf b}=\xb^{\mathbf c}\inn_\prec(g)$.
\TODO{form $f-\xb^{\mathbf c}g$: it lies in $I$, has strictly smaller initial monomial, contradiction}
\end{proof}

\begin{corollary}[Hilbert Basis Theorem, Cor.\ 1.14]\label{cor:hbt}
Every ideal in $K[\xb]$ is \fillin[3.5cm].
\end{corollary}

\block{Reduced Gröbner bases and standard monomials}

\begin{definition}[reduced, Def.\ 1.15]
$\mathcal G$ is a \keyword{reduced} Gröbner basis if (a) each leading coefficient is $1$, and (b) for
distinct $g,h\in\mathcal G$, no monomial of $g$ is a multiple of $\inn_\prec(h)$.
\end{definition}

\begin{theorem}[uniqueness + a basis, Thm.\ 1.16--1.17]\label{thm:reduced}
Fix $I$ and $\prec$. Then $I$ has a \emph{unique} reduced Gröbner basis. Moreover the
\keyword{standard monomials} $\mathcal S_\prec(I)=\{$monomials $\notin\inn_\prec(I)\}$ form a
$K$-vector-space basis of the quotient $K[\xb]/I$.
\end{theorem}
\begin{proof}[Proof of the basis claim]
Independence: a nonzero $f$ with image $0$ in $K[\xb]/I$ lies in $I$, so $\inn_\prec(f)\in\inn_\prec(I)$
is \emph{not} standard. Spanning: \TODO{take a minimal $\xb^{\mathbf c}$ not in the span; it lies in $\inn_\prec(I)$, divide, contradiction}
\end{proof}

\begin{yourturn}
(Symmetric polynomials, from Example 1.18.) The elementary symmetric polynomials
$\mathcal F=\{x+y+z,\ xy+xz+yz,\ xyz\}$ in $\Q[x,y,z]$ have reduced Gröbner basis (lex)
$\mathcal G=\{\underline{x}+y+z,\ \underline{y^2}+yz+z^2,\ \underline{z^3}\}$.
\begin{itemize}[leftmargin=1.3em,itemsep=0pt]
  \item List the \fillin[1cm] standard monomials $\mathcal S_\prec(I)$ (they index a basis of $\Q[x,y,z]/I$).
  \item This quotient is the regular representation of which group? \fillin[2cm]
\end{itemize}
\recall{\# standard monomials $=\dim_K K[\xb]/I=$ \# solutions counted with multiplicity. Here $=|S_3|=6$.}
\end{yourturn}

\begin{strand}
Buchberger's algorithm computes the reduced Gröbner basis from any $\mathcal F$. Two warnings the
book stresses: the \keyword{choice of order matters enormously} (Example 1.19: two random quartics in
$\PP^3$ give $5$ generators of degree $\le7$ in degrevlex, but $150$ of degree $\le73$ in lex), and
\keyword{lex is the order for elimination} (Chapter 4). Run a computer algebra system as you read.
\end{strand}

\loose{Exercise 1 (book): an ideal in $K[\xb]$ is principal $\iff$ its reduced Gröbner basis is a singleton. (One direction is easy; the other uses uniqueness.)}

% ===========================================================================
\section{Dimension and Degree}\label{sec:dimdeg}
\warmth{0}\quad\whisper{Two numbers attached to an ideal, both squeezed out of one counting series.}

\begin{strand}
The two most important invariants of an ideal are its \keyword{dimension} and \keyword{degree}.
Both are read off a single generating function that \emph{counts monomials not in the ideal}, the
Hilbert series. The definition here is purely combinatorial; §\ref{sec:affine} reveals that it
matches the geometric dimension and degree of the variety.
\end{strand}

\block{Counting what is left}

\begin{definition}[Hilbert function \& series, Def.\ 1.20--1.21]\label{def:hilbert}
For a monomial ideal $I\subset K[\xb]$, the \keyword{Hilbert function} $h_I(q)$ is the number of
monomials of degree $q$ that are \fillin[3cm] $I$. Its generating function is the
\keyword{Hilbert series} $\HS_I(z)=\sum_{q\ge0} h_I(q)\,z^q$.
\end{definition}

\begin{yourturn}
Warm up by counting monomials directly.
\begin{itemize}[leftmargin=1.3em,itemsep=1pt]
  \item Zero ideal $I=\{0\}$: every monomial survives, so $h_{\{0\}}(q)=\binom{n+q-1}{n-1}$ and
        $\HS_{\{0\}}(z)=\fillin[2cm]$. \recall{This $h$ is a polynomial in $q$ of degree $n-1$.}
  \item Principal $I=\gen{\xb^{\mathbf a}}$ of degree $e$: subtract the multiples, giving
        $\HS_I(z)=\dfrac{1-z^{e}}{(1-z)^{n}}$. Read off: $\dim=\fillin[1.2cm]$, $\deg=\fillin[1.2cm]$.
\end{itemize}
\end{yourturn}

\begin{theorem}[shape of the series, Thm.\ 1.25]\label{thm:hs}
For any monomial ideal $I\subset K[\xb]$,
\[
  \HS_I(z)\;=\;\frac{\kappa_I(z)}{(1-z)^{n}},\qquad \kappa_I\in\Z[z],\ \kappa_I(0)=1,
\]
and there is a \keyword{Hilbert polynomial} $\HP_I$ with $\HP_I(q)=h_I(q)$ for all large $q$.
\end{theorem}
\begin{proof}
Inclusion--exclusion over the minimal generators $m_1,\dots,m_r$. For $\tau\subseteq\{1,\dots,r\}$
let $e_\tau=\deg\,\mathrm{lcm}(m_i:i\in\tau)$. Then $\kappa_I(z)=\sum_\tau(-1)^{|\tau|}z^{e_\tau}$.
\TODO{regroup to get $h_I(q)=\sum_\tau(-1)^{|\tau|}\binom{n+q-e_\tau-1}{n-1}$; check it is a polynomial for $q\gg0$}
\end{proof}

\block{Reading off dimension and degree}

\begin{definition}[dimension \& degree, Def.\ 1.28]\label{def:dimdeg}
Write $\HP_I(q)=\dfrac{g}{(d-1)!}\,q^{\,d-1}+(\text{lower order})$. Then the \keyword{dimension} of
$I$ is $d$ and the \keyword{degree} is the positive integer $g$. If $\HP_I\equiv0$ then $\dim=0$ and
$\deg I=\dim_K K[\xb]/I$ (a finite number).
\end{definition}

\recall{Slogan: dimension is the \emph{growth rate} of $h_I$; degree is its \emph{leading coefficient} (times $(d-1)!$).}

\begin{yourturn}
Two computations the book recommends.
\begin{itemize}[leftmargin=1.3em,itemsep=1pt]
  \item (Ex.\ 1.31) Principal ideal generated by a single monomial of degree $e$:
        confirm $\dim=n-1$ and $\deg=e$ from $\HP_I(q)=\dfrac{e}{(n-2)!}q^{\,n-2}+\cdots$.
  \item (Ex.\ 1.32) $n=2m$, $I=\gen{x_1x_2,\,x_3x_4,\,\dots,\,x_{2m-1}x_{2m}}$:
        find $\dim I=\fillin[1.2cm]$ and $\deg I=\fillin[1.2cm]$. \recall{Try $m=2$: $I=\gen{x_1x_2,x_3x_4}$. The degree is a power of two.}
\end{itemize}
\workspace[2]
\end{yourturn}

\begin{strand}
Why these two numbers (geometry preview, Remark 1.36 / Example 1.37). For the variety $\V(I)$:
$\dim=0$ \emph{iff} it is finitely many points, and then $\deg=$ the \emph{number} of points (counted
by the radical). A curve has $\dim 1$, a surface $\dim 2$. The degree measures how curvy the shape is.
A clean fact: a prime ideal has degree $1$ \emph{iff} it is generated by linear polynomials. For the
running cubic, $\gen f$ has $\dim 2,\deg 3$; its singular locus has $\dim 0,\deg 5$; together with $f$,
$\dim 0,\deg 4$ (the four nodes).
\end{strand}

\loose{Exercise 14 (book): let $X$ be a generic $2\times2$ matrix of variables, $I_s=\gen{\text{entries of }X^s}$. Investigate $\dim,\deg$ of $I_s$ for $s=2,3,4,\dots$. A small experiment with a CAS pays off.}

% ===========================================================================
\section{Affine Varieties}\label{sec:affine}
\warmth{0}\quad\whisper{The geometry side, and the bridge. ``Geometry is but drawn algebra.''}

\begin{strand}
A \keyword{variety} is the solution set of polynomial equations,
$\V(I)=\{p\in K^n: f(p)=0\ \forall f\in I\}$. Different ideals can give the same variety, so geometry
sees only $I$ \emph{up to radical}. The reverse map $\Ideal(W)=\{f:f|_W=0\}$ sends a shape back to its
(radical) ideal. These two maps are the dictionary; this section makes them precise.
\end{strand}

\block{The two maps $\V$ and $\Ideal$}

\begin{definition}[variety and vanishing ideal]\label{def:VandI}
For an ideal $I\subseteq K[\xb]$ and a subset $W\subseteq K^n$,
\[
  \V(I)=\{p\in K^n: f(p)=0\ \forall f\in I\},\qquad
  \Ideal(W)=\{f\in K[\xb]: f(p)=0\ \forall p\in W\}.
\]
$\Ideal(W)$ is always a \keyword{radical} ideal. A subset is a \emph{variety} iff $W=\V(\Ideal(W))$.
\end{definition}

\recall{Both maps reverse inclusions. Fill: $I\subseteq J\Rightarrow \V(I)\,\fillin[1cm]\,\V(J)$, and $V\subseteq W\Rightarrow\Ideal(W)\,\fillin[1cm]\,\Ideal(V)$.}

\block{Irreducible $=$ prime: the keystone}

\begin{definition}[irreducible]
A variety is \keyword{irreducible} if it is not a finite union of proper subvarieties:
$\V(I)=\V(J)\cup\V(J')\Rightarrow\V(I)=\V(J)$ or $\V(J')$.
\end{definition}

\begin{proposition}[Prop.\ 2.3]\label{prop:irr-prime}
A variety $W\subseteq K^n$ is irreducible \;$\iff$\; its ideal $\Ideal(W)$ is \fillin[2.5cm].
\end{proposition}
\begin{proof}
($\Leftarrow$) Suppose $\Ideal(W)$ prime and $W=\V(J)\cup\V(J')$ with $W\ne\V(J)$. Pick $f\in J$,
$v\in W$ with $f(v)\ne0$, so $f\notin\Ideal(W)$. For $g\in J'$, the product $fg$ vanishes on
$\V(J)\cup\V(J')=W$, so $fg\in\Ideal(W)$; primality gives $g\in\Ideal(W)$. Hence $W\subseteq\V(J')$.
\TODO{($\Rightarrow$) $W$ irreducible: if $fg\in\Ideal(W)$ then $W=(W\cap\V(f))\cup(W\cap\V(g))$; conclude $f$ or $g\in\Ideal(W)$}
\end{proof}

\begin{strand}
This is the dictionary's keystone: \keyword{prime ideals} $\leftrightarrow$ \keyword{irreducible
varieties}, exactly as prime numbers are to $\Z$. In applications (algebraic statistics), a model is
the image of a polynomial map; its Zariski closure is irreducible, so its ideal is prime, and is
computed as the kernel of the dual ring map. The closed sets of the \keyword{Zariski topology} are
precisely the varieties; the coordinate ring $K[W]=K[\xb]/\Ideal(W)$ is the ring of polynomial
functions on $W$, and a map of varieties $f:W_1\to W_2$ pulls back to a ring map
$f^\ast:K[W_2]\to K[W_1]$ (a \emph{contravariant} functor).
\end{strand}

\block{Dimension and degree, geometrically}

\begin{example}[linear subspace, Ex.\ 2.15]
For a linear subspace $V\subseteq K^n$: $\dim V$ is its linear-algebra dimension, $\Ideal(V)=
\gen{x_1,\dots,x_s}$ with $s=n-\dim V$, and $\deg V=\fillin[1cm]$. \recall{Degree $1$ $\iff$ prime ideal generated by linear forms.}
\end{example}

A method to get $\dim\V(I)$: compute $\inn_\prec(I)$; find the smallest variable set $S$, $|S|=d$, hitting
every generator; then $\dim=n-d$.

\begin{theorem}[Bézout, Thm.\ 2.16]\label{thm:bezout}
Let $f_1,\dots,f_k$ be \emph{general} polynomials of degrees $d_1,\dots,d_k>0$ in $n$ variables, and
$I=\gen{f_1,\dots,f_k}$. Then
\[
  \dim(I)=\fillin[1.5cm],\qquad \deg(I)=\fillin[2cm].
\]
(Such $I$ is a \keyword{complete intersection}; $\dim\ge n-k$ always, by Krull.)
\end{theorem}
\begin{yourturn}
Two general conics in the plane ($n=2,k=2,d_1=d_2=2$): Bézout predicts $\dim=\fillin[0.8cm]$ and
$\deg=\fillin[0.8cm]$ -- i.e.\ they meet in $\fillin[0.8cm]$ points (over $\overline K$). Sketch two
ellipses meeting in $4$ points to see it.
\end{yourturn}

\block{Smooth vs singular}

\begin{definition}[singular point]
$p\in\V(f)$ is \keyword{singular} if all partials vanish: $\partial f/\partial x_i(p)=0$ for all $i$.
The singular locus is $\V\!\big(\gen{f,\partial f/\partial x_1,\dots,\partial f/\partial x_n}\big)$; if
empty, $\V(f)$ is \keyword{smooth}. For an irreducible $Y$ of dimension $d$ cut by a prime $I$, the point
$p$ is singular iff $\rank\,\mathrm{Jac}(p)<n-d$ (codimension).
\end{definition}

\recall{The four nodes of the running cubic are exactly its singular locus -- the $\dim 0,\deg 4$ ideal from §\ref{sec:dimdeg}. Smoothness $=$ the tangent space has the right dimension.}

\loose{Bertini (book): an ideal generated by general polynomials is prime and its variety is smooth, provided $\dim>0$. So ``random'' equations cut out nice irreducible smooth varieties; degeneracy is special.}

% ===========================================================================
\section{Projective Varieties}\label{sec:proj}
\warmth{0}\quad\whisper{Add points at infinity and the geometry gets simpler: parallel lines finally meet.}

\begin{strand}
Projective space $\PP^n=\PP(K^{n+1})$ is the set of \emph{lines through the origin} in $K^{n+1}$,
with coordinates $[a_0:\cdots:a_n]$ defined up to scaling. It is covered by $n+1$ affine charts
$\{a_i\ne0\}\cong K^n$, and over $\R,\C$ it is \keyword{compact}. Geometers prefer $\PP^n$: parallel
lines in $K^2$ do not meet, but any two lines in $\PP^2$ do. The price is that we may only use
\keyword{homogeneous} polynomials, since only their vanishing is well defined on a line.
\end{strand}

\block{Homogeneous ideals, cones, closures}

\begin{definition}[projective variety]
For homogeneous $f_1,\dots,f_k$, $\ \V(f_1,\dots,f_k)=\{[a]\in\PP^n: f_i(a)=0\ \forall i\}$. An ideal is
\keyword{homogeneous} if generated by homogeneous polynomials. The \keyword{affine cone} $\hat X\subseteq
K^{n+1}$ is cut by the same ideal, and
\[
  \dim(X)=\dim(\hat X)-\fillin[0.8cm],\qquad \deg(X)=\deg(\hat X).
\]
\end{definition}

\recall{Why is the cone $\hat X$ singular at the origin once $\deg\ge2$? Every line of the cone passes through $0$.}

The \keyword{projective closure} $\bar Y\subseteq\PP^n$ of an affine $Y\subseteq K^n$ has ideal $\bar I$
obtained by \keyword{homogenizing}, and Proposition 2.18 computes it cleanly: take a degree-compatible
reduced Gröbner basis $\mathcal G$ of $I$ and homogenize each $g\in\mathcal G$.

\begin{theorem}[invariants survive the closure, Cor.\ 2.19]\label{cor:closure}
$\dim(\bar Y)=\dim(Y)$ and $\deg(\bar Y)=\deg(Y)$.
\end{theorem}

\begin{yourturn}
(Rational normal curve, Ex.\ 2.20.) Let $I=\gen{x_i-x_1^i:i=2,\dots,n}$, so $Y=\V(I)$ is the curve
$t\mapsto(t,t^2,\dots,t^n)$. Its projective closure $\bar Y\subset\PP^n$ has ideal generated by the
$2\times2$ minors of
\[
  \begin{pmatrix} x_0 & x_1 & x_2 & \cdots & x_{n-1}\\ x_1 & x_2 & x_3 & \cdots & x_n \end{pmatrix}.
\]
How many such minors are there? $\fillin[1.5cm]$. The curve $\bar Y$ has degree $\fillin[1cm]$ and is
the \keyword{rational normal curve}. \recall{The $\binom n2$ minors; degree $n$. The initial ideal is $\gen{x_1,\dots,x_{n-1}}^2$.}
\end{yourturn}

\block{Why projective: intersections always behave}

\begin{theorem}[projective intersection, Thm.\ 2.22]\label{thm:intersect}
Let $K$ be algebraically closed and $X,Y\subseteq\PP^n$ projective varieties with $d_1=\dim X$,
$d_2=\dim Y$. Then
\[
  \dim(X\cap Y)\;\ge\; \fillin[2.5cm].
\]
In particular, if $d_1+d_2\ge n$ then $X\cap Y$ is \fillin[2.5cm].
\end{theorem}

\begin{yourturn}
(Ex.\ 2.23, why the hypotheses matter.) Two quadric surfaces $X,Y\subseteq\PP^3$ ($n=3$, $d_1=d_2=2$).
\begin{itemize}[leftmargin=1.3em,itemsep=1pt]
  \item Theorem prediction: $\dim(X\cap Y)\ge 2+2-3=\fillin[0.8cm]$, and the intersection is nonempty.
  \item Over $\C$: their intersection is a union of four lines, $\dim=1$. \checkmark
  \item Over $\R$: the book's pair meets in only \emph{two points}, $\dim=0$. Which hypothesis of the
        theorem just failed? \fillin[3cm]
\end{itemize}
\end{yourturn}

\begin{strand}
Many models in science are projective varieties given by homogeneous constraints. Two showcases:
the \keyword{nilpotent $n\times n$ matrices} form an irreducible projective variety in $\PP^{n^2-1}$ of
$\dim n^2-n-1$ and $\deg n!$, a complete intersection cut by the non-leading coefficients of the
characteristic polynomial (for $n=2$: $\Ideal(X)=\gen{\operatorname{trace}(A),\det(A)}$); and the
\keyword{Kalman varieties} of matrices with an eigenvector in a fixed subspace, basic in control theory.
The lesson of the chapter: when constraints are homogeneous, regard the model as living in $\PP^n$.
\end{strand}

\loose{Exercise 17 (book): $X\subset\C^4$ cut by two random cubics. What does Bézout tell you about $\dim,\deg$, and what does Bertini tell you about smoothness and primeness?}

% ===========================================================================
\section{Geometry in Low Dimensions}\label{sec:lowdim}
\warmth{0}\quad\whisper{Where equations become manifolds: circles, Riemann spheres, ovals, and a torus.}

\begin{strand}
Smooth projective varieties in low dimensions \emph{are} familiar manifolds. Over $\R$ a smooth
projective variety of dimension $d$ is a compact real manifold of dimension $d$; over $\C$, of
dimension $2d$. So $\PP^1_{\R}$ is a \fillin[2cm] and $\PP^1_{\C}$ is the \keyword{Riemann sphere}. This
section reads the topology straight off the algebra (degree, genus).
\end{strand}

\block{Curves in the plane: ovals, pseudolines, genus}

\begin{yourturn}
($n=1$, binary forms.) A proper subvariety of $\PP^1_{\C}$ is the zero set of a binary form
$f(x,y)$ (homogeneous in two variables). For $f=x^{11}y-11x^6y^6-xy^{11}$, the variety $\V(f)$ has
$\dim 0$ and degree $\fillin[1cm]$. These $12$ points on the Riemann sphere are Klein's
\keyword{icosahedron} vertices: $4$ real, $8$ in conjugate pairs. \recall{Degree of a binary form $=$ its number of roots in $\PP^1_{\C}$, with multiplicity.}
\end{yourturn}

For a curve $C\subset\PP^2_{\R}$, a connected component is a \keyword{pseudoline} if removing it leaves
$\PP^2_{\R}$ connected, and an \keyword{oval} otherwise (it has an inside and an outside, like a conic).

\begin{theorem}[curves in the plane, Thm.\ 2.26]\label{thm:curve}
Let $C$ be a smooth curve of degree $d$ in $\PP^2_K$.
\begin{itemize}[leftmargin=1.3em,itemsep=1pt]
  \item If $K=\C$: $C$ is an orientable surface of \keyword{genus} $g=\dfrac{(d-1)(d-2)}{2}$.
  \item If $K=\R$: $C$ has at most $g+1$ components; if $d$ is even all are ovals, if $d$ is odd exactly
        one is a \fillin[2.5cm].
\end{itemize}
\end{theorem}
\begin{yourturn}
Fill the genus by degree: $d=1\Rightarrow g=\fillin[0.6cm]$, $\ d=2\Rightarrow g=\fillin[0.6cm]$,
$\ d=3\Rightarrow g=\fillin[0.6cm]$. So lines and conics are spheres ($g=0$), and a smooth cubic is a
\fillin[1.5cm] ($g=1$).
\end{yourturn}

\block{Elliptic curves: a cubic that is a torus}

\begin{strand}
A smooth cubic ($d=3,g=1$), e.g.\ the Weierstrass curve $zy^2=x^3-xz^2$, is an \keyword{elliptic
curve}. Over $\R$ it shows an oval and a pseudoline; over $\C$ it is homeomorphic to a
\keyword{torus} $S^1\times S^1=\R^2/M$, where $M=\Z\lambda+\Z\tau$ is the period lattice obtained by
integrating $\int \mathrm dx/y$ along loops. The torus carries a group law, the seed of $19$th-century
geometry and modern elliptic-curve cryptography. \recall{Beware two ``tori'': this \emph{topological} torus is \emph{not} the algebraic torus $(\C^\ast)^n$ of toric geometry (Ch.\,8). An elliptic curve is not toric.}
\end{strand}

\block{The dictionary, completed}

\begin{strand}
Look back at the front-page table and check every row you can now fill. The two chapters were one
conversation, in two languages:
\end{strand}
\noindent\begin{tabularx}{\linewidth}{@{}L L@{}}
\toprule
{\color{indigo}Algebra} & {\color{madder}Geometry}\\
\midrule
prime ideal & irreducible variety \;(Prop.\ \ref{prop:irr-prime})\\
maximal ideal & a point\\
$I\subseteq J$ & $\V(I)\supseteq\V(J)$ \;(inclusion-reversing)\\
$I+J$ \;/\; $I\cap J$ & intersection \;/\; union\\
dimension, degree of $I$ & dimension, degree of $\V(I)$ \;(they agree!)\\
homogeneous ideal & projective variety\\
$\Ideal(\V(I))=\rad(I)$ & the Nullstellensatz bridge \;(Ch.\,6)\\
primary decomposition & decomposition into irreducibles \;(Ch.\,3)\\
\bottomrule
\end{tabularx}

\begin{strand}
Where this goes next. Chapter 3 (\emph{Solving and Decomposing}) turns ``$\V(I)=$ union of
irreducibles'' into an algorithm: \keyword{primary decomposition}, the geometry behind the Chinese
Remainder Theorem you already met in the running cubic. Keep the cubic, the dictionary, and the
Gröbner engine in hand; everything later is built from these three.
\end{strand}

\loose{Exercise 15 (book): solve Diophantus' problem -- find a rational point on the elliptic curve $y(6-y)=x^3-x$. Hint: take the tangent line at $(-1,0)$ and reintersect.}


\end{document}
